%%%%%%%%%%%%%%%%%%%%%%%%%%%%%%%%%%%%%%%%%%%%%%%%%%%%%%%%%%%%%%%%%%%%%%
% Lớp 10 | Đại số tổ hợp | Hoán vị. Chỉnh hợp. Tổ hợp
% Dạng 0D8?2-1: Lý thuyết tổng hợp về 𝑃, 𝐶, 𝐴
% Số câu: 18  (sinh tự động bởi tach_cau.py; từ nay thầy sửa trực tiếp file này)
%%%%%%%%%%%%%%%%%%%%%%%%%%%%%%%%%%%%%%%%%%%%%%%%%%%%%%%%%%%%%%%%%%%%%%

%%%%%%%%%%%%%%%% TRẮC NGHIỆM 4 PHƯƠNG ÁN %%%%%%%%%%%%%%%%

% nguồn: MaHoaNH_TracNghiem_Lop10
\begin{ex}%[0D8N2-1]%[ma:3c50827a]
	Cho $n\in{\mathbb{N}^{*}}$ thỏa mãn $\mathrm{C}_n^5=2002$. Tính $\mathrm{A}_n^5$. 
	 \choice
		{\True $240\  240$}
		{$10\ 010$} 
		{$40\ 040$} 
		 {$2\ 007$}
	\loigiai{ 
	Ta có $\mathrm{A}_n^5=\mathrm{C}_n^5\cdot  5!=240\ 240$.}
\end{ex}

% nguồn: MaHoaNH_TracNghiem_Lop10
\begin{ex}%[0D8N2-1]%[ma:869c5605]
	 Có bao nhiêu cách sắp xếp $5$ học sinh thành một hàng dọc?
	\choice 
		 {$5$}
		 {\True $120$}
		{$5^5$} 
		 {$24$}
	 \loigiai{
		Sắp $5$ học sinh vào $5$ vị trí hàng dọc có $5!$ cách. 
	 }
\end{ex}

% nguồn: MaHoaNH_TracNghiem_Lop10
\begin{ex}%[0D8N2-1]%[ma:6cdcb1de]
	Số các  hoán vị của $7$ phần tử là
	 \choice
		{$7$} 
		 {$6$}
		{\True  $7!$}
		{$6!$} 
	 \loigiai{
		Số các hoán vị của $7$ phần tử là $7!$. 
	} 
\end{ex}

% nguồn: MaHoaNH_TracNghiem_Lop10
\begin{ex}%[0D8N2-1]%[ma:57295c27]
	Cho $n$ là một số nguyên dương $n \geq 9$. Số các chỉnh hợp  chập $9$ của $n$ phần tử là
	\choice 
		{$\mathrm{C}_n^9$} 
		{$9!$} 
		{$n!$} 
		{\True  $\mathrm{A}_n^9$}
	 \loigiai{
		Số các chỉnh  hợp chập $9$ của $n$ phần tử là $\mathrm{A}_n^9$.
	 }
\end{ex}

% nguồn: MaHoaNH_TracNghiem_Lop10
\begin{ex}%[0D8N2-1]%[ma:9e65bc53]
	Số các tổ hợp chập $6$ của  $9$ phần tử là
	 \choice
		 {$\mathrm{C}_9^6$}
		 {$6!$}
		 {\True $\mathrm{A}_9^6$}
		{$9!$} 
	 \loigiai{
		Số các  tổ hợp chập $6$ của $9$ phần tử là $\mathrm{A}_9^6$.
	 }
\end{ex}

% nguồn: MaHoaNH_TracNghiem_Lop10
\begin{ex}%[0D8N2-1]%[ma:48168471]
	Số tổ hợp chập $k$ của $n$ phần tử ($k\le n$)  là
	 \choice
		 {$\mathrm{C}_n^k=\dfrac{n}{k!\left( n-k \right)!}$}
		{$\mathrm{C}_n^k=\dfrac{n!}{k\left(  n-k \right)!}$}
		{$\mathrm{C}_n^k=\dfrac{n}{k!\left(  n+k \right)!}$}
		{\True $\mathrm{C}_n^k=\dfrac{n!}{k!\left( n-k \right)!}$} 
	 \loigiai{
		Số  tổ hợp chập $k$ của $n$ phần tử ($k\le n$) là $\mathrm{C}_n^k=\dfrac{n!}{k!\left( n-k \right)!}$.
	} 
\end{ex}

% nguồn: MaHoaNH_TracNghiem_Lop10
\begin{ex}%[0D8N2-1]%[ma:4a52a2da]
	Có bao nhiêu  cách xếp $6$ học sinh thành một hàng dọc?
	 \choice
		 {$\mathrm{C}_6^2$}
		 {$\mathrm{A}_6^2$}
		{\True $6 !$} 
		{$6^6$} 
	\loigiai{ 
		Số cách sắp xếp $6$ học sinh thành  một hàng dọc là số hoán vị của $6$ học sinh.\\
		 Vậy có tất cả $6!$ cách sắp xếp.
	 }
\end{ex}

% nguồn: MaHoaNH_TracNghiem_Lop10
\begin{ex}%[0D8H2-1]%[ma:e365acac]
	Có bao nhiêu cách chọn $2$ học sinh từ một tổ  gồm có $9$ học sinh giữ chức danh tổ trưởng và tổ phó?
	 \choice
		{$2^9$} 
		{$\mathrm{C}_9^2$} 
		{$9^2$} 
		{\True $\mathrm{A}_9^2$} 
	\loigiai{ 
		Số cách chọn $2$ học sinh từ một tổ gồm có $9$ học sinh giữ chức danh  tổ trưởng và tổ phó là $\mathrm{A}_9^2$.
	} 
\end{ex}

% nguồn: MaHoaNH_TracNghiem_Lop10
\begin{ex}%[0D8H2-1]%[ma:20720b8a]
	Số các hoán  vị của $4$ phần tử là
	\choice 
		{\True $24$} 
		{$12$} 
		{$4$} 
		 {$48$}
	\loigiai{ 
		Áp dụng công thức $\mathrm{P}_n=n!\Rightarrow  \mathrm{P}_4=24$.
	} 
\end{ex}

% nguồn: MaHoaNH_TracNghiem_Lop10
\begin{ex}%[0D8H2-1]%[ma:748dfb64]
	Có  bao nhiêu cách xếp $8$ học sinh thành một hàng ngang?
	\choice 
		{\True $8!$} 
		 {$1$}
		 {$8$}
		 {$8^8$}
	 \loigiai{
		Mỗi cách xếp $8$ em học sinh thành một hàng ngang  là một hoán vị của $8$ phần tử. Vậy có $8!$ cách xếp $8$ học sinh thành một hàng ngang.
	 }
\end{ex}

% nguồn: MaHoaNH_TracNghiem_Lop10
\begin{ex}%[0D8H2-1]%[ma:3747692b]
	Từ tập $A=\{1;2;3;4;5;6;7;8;9\}$ lập được  bao nhiêu số tự nhiên chẵn có $4$ chữ số khác nhau?
	 \choice
		{$A^4_9$} 
		 {\True $4A^3_8$}
		{$4A^3_9$} 
		 {$9A^3_8$}
	 \loigiai{
		Có $4$ cách chọn chữ số hàng đơn vị sao cho chữ số hàng đơn  vị là số chẵn.\\
		Ba chữ số  ở các hàng còn lại ta chọn $3$ trong $8$ số của tập $A$ (đã trừ đi $1$ số chẵn đứng ở hàng đơn vị).\\
		Vậy có tất cả $4A^3_8$ số thỏa mãn  đề bài.
	 }
\end{ex}


%%%%%%%%%%%%%%%% ĐÚNG/SAI %%%%%%%%%%%%%%%%

% nguồn: MaHoaNH_DungSai_Lop10
\begin{ex}%[0D8N2-1]%[ma:80f67b57]
	Số các tổ hợp  chập $k$ của $n$ phần tử là
	 \choiceTF
		{\True  $\mathrm{C}_n^k=\dfrac{n!}{(n-k)!(k-1)!\cdot k}$}
		 {$\mathrm{A}_n^k=\dfrac{n!}{(n-k)!k!}$}
		{\True  $\mathrm{C}_n^k=\dfrac{n!}{(n-k)!k!}$}
		 {$\mathrm{C}_n^k=\dfrac{n!}{(n-k)!}$}
	 \loigiai{
	Số  các tổ hợp chập $k$ của $n$ phần tử là $\mathrm{C}_n^k=\dfrac{n!}{(n-k)!k!}==\dfrac{n!}{(n-k)!(k-1)!\cdot k}$.}
\end{ex}

% nguồn: MaHoaNH_DungSai_Lop10
\begin{ex}%[0D8N2-1]%[ma:9f9ba330]
	Số các chỉnh  hợp chập $k$ của $n$ phần tử là
	 \choiceTF
		 {\True $\mathrm{A}_n^k=\dfrac{n!}{(n-k)!}$}
		{\True $\mathrm{A}_n^k=\dfrac{(n-1)!\cdot  n}{(n-k)!}$}
		 {$\mathrm{C}_n^k=\dfrac{n!}{(n-k)!k!}$}
		 {$\mathrm{C}_n^k=\dfrac{n!}{(n-k)!}$}
	\loigiai{ 
	Số các  chỉnh hợp chập $k$ của $n$ phần tử là $\mathrm{A}_n^k=\dfrac{n!}{(n-k)!}=\dfrac{(n-1)!\cdot n}{(n-k)!}$.}
\end{ex}

% nguồn: MaHoaNH_DungSai_Lop10
\begin{ex}%[0D8N2-1]%[ma:ab3b46f3]
	Xét tính đúng  sai của các mệnh đề sau:
	\choiceTF{ Cho tập $ A $ gồm $ n $ phần tử và $ 1\le k\le n $. Kết quả của việc lấy ra $ k $ phần tử từ $ n $ phần tử từ tập hợp $ A $ và sắp xếp chúng theo một thứ tự nào  đó được gọi là một tổ hợp chập $ k $ của $ n $ phần tử đã cho }
		 {\True Với $ n $ là số nguyên dương bất kì $ n\ge3 $ thì ta có $ \mathrm{A}_n^3=\dfrac{n!}{(n-3)!} $}
		{\True Số chỉnh hợp chập $ k $ của $ n $  phần tử gấp $ k! $ lần số tổ hợp chập $ k $ của $ n $ phần tử}
		{\True Với $  n $ nguyên dương bất kì và $ n\ge3 $ thì ta có $ \mathrm{C}_n^3=\dfrac{n!}{3!(n-3)!} $}
	\loigiai{ 
		\begin{itemchoice}
			\itemch Cho tập $ A $ gồm $ n $ phần tử và $ 1\le k\le n $. Kết  quả của việc lấy ra $ k $ phần tử từ $ n $ phần tử từ tập hợp $ A $ và sắp xếp chúng theo một thứ tự nào đó được gọi là một \textbf{chỉnh hợp} chập $ k $ của $ n $ phần tử đã cho.
			\itemch Với $ n $  là số nguyên dương bất kì $ n\ge3 $ thì ta có $ \mathrm{A}_n^3=\dfrac{n!}{(n-3)!}$.
			\itemch  Ta có $ \mathrm{A}_n^k=k! \mathrm{C}_n^k$.
			\itemch Ta có $  \mathrm{C}_n^3=\dfrac{n!}{3!(n-3)!}$.
		\end{itemchoice}
	} 
\end{ex}


%%%%%%%%%%%%%%%% TỰ LUẬN %%%%%%%%%%%%%%%%

% nguồn: MaHoaNH_TuLuan_Lop10
\begin{bt}%[0D8Y2-1]%[ma:9e77c9f4]
	Tổ $ 1 $ có $ 12 $ học sinh trong đó có $ 4 $ học sinh nam và $ 8 $ học sinh nữ. Có bao nhiêu cách chọn $ 3 $ học sinh  làm lớp trưởng, lớp phó và thủ quỹ. (Mỗi bạn chỉ làm $ 1 $ nhiệm vụ)
	\loigiai{ 
	Mỗi cách chọn $ 3 $ học sinh từ $ 12 $ học sinh xếp vào các vị trí lớp trưởng, lớp phó và thủ quỹ là một chỉnh hợp chập $ 3 $ của $ 12 $ phần tử.\\ 
	Suy ra có $ \mathrm{A}_{12}^3=1320 $  cách chọn.
	 }
\end{bt}

% nguồn: MaHoaNH_TuLuan_Lop10
\begin{bt}%[0D8B2-1]%[ma:de9ee121]
	Trong một lô $100$ sản phẩm, có $96$ chính phẩm (sản phẩm đạt tiêu chuẩn) và $4$  thứ phẩm (sản phẩm không đạt tiêu chuẩn). Từ $100$ sản phẩm này có bao nhiêu cách lấy ra $3$ sản phẩm mà trong đó có ít nhất một thứ phẩm?
	\loigiai{ 
		Chọn $3$ sản phẩm  bất kì có $\mathrm{C}^3_{100}=161700$ cách.\\
		Chọn $3$ sản  phẩm mà không có thứ phẩm nào có $\mathrm{C}^3_{96}=142880$ cách.\\
		Suy ra số cách lấy ra $3$  sản phẩm mà trong đó có ít nhất một thứ phẩm là
		$$161700 -  142880 =18820 \text{ (cách).}$$
	} 
\end{bt}

% nguồn: MaHoaNH_TuLuan_Lop10
\begin{bt}%[0D8B2-1]%[ma:8b9c07cd]
	Một tổ có $12$ học sinh. Hỏi có bao nhiêu cách chia $12$ học sinh ra làm $4$  nhóm $A$, $B$, $C$ và $D$ sao cho mỗi nhóm có số học sinh bằng nhau?
	\loigiai{ 
		Do có $12$ học sinh chia  đều ra $4$ nhóm với số học sinh bằng nhau nên mỗi nhóm có $3$ học sinh.\\
		Chọn $3$ học sinh vào  nhóm $A$ có $\mathrm{C}_{12}^3$ cách.\\
		Chọn $3$ học sinh  vào nhóm $B$ có $\mathrm{C}_9^3$ cách.\\
		Chọn $3$  học sinh vào nhóm $C$ có $\mathrm{C}_6^3$ cách.\\
		Chọn $3$ học sinh vào  nhóm $D$ có $\mathrm{C}_3^3$ cách.\\
		Vậy  có $\mathrm{C}_{12}^3\cdot\mathrm{C}_9^3\cdot\mathrm{C}_6^3\cdot\mathrm{C}_3^3=369600$ cách.
	 }
\end{bt}

% nguồn: MaHoaNH_TuLuan_Lop10
\begin{bt}%[0D8Y2-1]%[ma:05a5361c]
	Từ các chữ số $0$; $1$; $2$; $3$; $4$; $5$; $6$  có thể lập được bao nhiêu số tự nhiên lẻ có năm chữ số đôi một khác nhau?
	 \loigiai{
		 Gọi số cần tìm là $\overline{abcde}$.
		\begin{itemize}
			\item Chọn  $e\in \{1;3;5\}$: có $3$ cách.
			\item Chọn $a\neq  0, a\neq e$: có $5$ cách.
			\item Chọn $3$ số từ $5$ số còn lại xếp vào vị trí $b$, $c$,  $d$: có $\mathrm{A}^3_5=60$ cách.
		\end{itemize}
		Vậy có tất cả $3\cdot 5\cdot 60 = 900$  số.
	 }
\end{bt}
